\documentclass[aspectratio=169,11pt]{beamer}
\usepackage[utf8]{inputenc}
\usepackage[T1]{fontenc}
\usepackage[english]{babel}
\usepackage{amsmath,amssymb}
\usepackage{graphicx}
\usetheme{default}
\usecolortheme{seahorse}
\setbeamertemplate{navigation symbols}{}
\setbeamertemplate{footline}[frame number]
\setbeamerfont{frametitle}{size=\Large}
\setbeamerfont{title}{size=\LARGE}
\setbeamersize{text margin left=7mm,text margin right=7mm}

\definecolor{accent}{RGB}{55,76,118}
\setbeamercolor{structure}{fg=accent}
\setbeamercolor{alerted text}{fg=red!70!black}

\title{The Economics of Bicycles for the Mind}
\subtitle{Agrawal, Gans \& Goldfarb (2025)\\[0.4em]
NBER Working Paper 34034}
\author{}
\date{}

\begin{document}

% 1 ------------------------------------------------------------------
\begin{frame}[plain]
  \titlepage
  \vspace{-1.1em}
  \begin{center}
    \small\url{https://github.com/fgalarzac/ai-02-agrawal}

    \vfill
    \footnotesize Repository 2 --- Artificial Intelligence and Economic Modeling
  \end{center}
\end{frame}

% 2 ------------------------------------------------------------------
\begin{frame}{The Paper: Cognitive Tools and Human Capabilities}
  \textbf{Question.} Why does the empirical record on cognitive technology and
  inequality look contradictory --- computers raised productivity \emph{and}
  inequality, while AI often raises productivity while \emph{reducing} it
  within firms?

  \smallskip
  \textbf{The single mechanism.} A cognitive tool is a substitute for
  implementation skill $s$, but a complement to judgment --- and judgment
  splits into two skills that behave differently:
  \begin{columns}[T,onlytextwidth]
    \column{0.48\textwidth}
      \textbf{Opportunity judgment $\gamma$}
      \begin{itemize}
        \item Spotting a chance to improve the task
        \item \emph{Always} complemented by a better tool
      \end{itemize}
    \column{0.48\textwidth}
      \textbf{Payoff judgment $\alpha$}
      \begin{itemize}
        \item Acting correctly on a successful implementation
        \item Complemented only \emph{conditionally}
      \end{itemize}
  \end{columns}

  \vspace{-0.1em}
  \begin{block}{Agent's problem}
    \vspace{-0.7em}
    \[
      \max_{e_t\geq0}\; M(e_t;\theta)=p(se_t;\theta)\,\alpha\,\Delta-c(e_t;\theta)
    \]
    \vspace{-1.0em}
    \[
      \text{FOC:}\quad p'(se_t^*;\theta)\,s\alpha\Delta=c'(e_t^*;\theta)
    \]
  \end{block}
\end{frame}

% 3 ------------------------------------------------------------------
\begin{frame}{The Main Result: Proposition 1, with All Its Conditions}
  \small
  \begin{block}{Definition 1 (cognitive tool)}
    \vspace{-0.6em}
    For all $e$ and all $\theta'>\theta$:
    \[
      \text{(i) } p(se;\theta')\geq p(se;\theta),\ \ c(e;\theta')\leq c(e;\theta)
      \qquad\qquad
      \text{(ii) } \frac{p'(se;\theta)}{c'(e;\theta)}\ \text{strictly decreasing in }\theta
    \]
  \end{block}

  \vspace{-0.3em}
  \begin{block}{Proposition 1}
    \vspace{-0.5em}
    Under Definition 1 and the standing regularity assumptions ($p$ increasing,
    weakly concave; $c$ increasing, weakly convex), as $\theta$ rises from $0$
    to $1$:
    \[
      \begin{gathered}
        \text{1. } e_t^*(1)<e_t^*(0)\ \text{for all }t
        \qquad\qquad
        \text{2. } e_t^*(\theta)=e^*(\theta)\ \text{for all }t\\[0.3em]
        \text{3. } V_0(1)>V_0(0)
      \end{gathered}
    \]
  \end{block}

  \vspace{-0.2em}
  \centering\emph{A better tool lowers the marginal payoff of effort relative
  to its cost, so the agent works less per opportunity --- but net benefit
  still rises, so does the value of the whole process.}
\end{frame}

% 4 ------------------------------------------------------------------
\begin{frame}{What I Did}
  \small
  \textbf{1. Audited the variance-U-shape claim (Prop.\ 3(b), Appendix A.3).}
  The naive claim --- ``variance is U-shaped in $\theta$, no conditions'' ---
  fails: it holds only under condition (30). But the paper's own proof text
  goes further: it states ``therefore $B<0$'' as if condition (30) is what
  establishes $B<0$. In fact $B<0$ follows from Jensen's inequality alone,
  \emph{independent} of condition (30); the real conditional fact is whether
  the $\Gamma$-heterogeneity term can overturn it in the full expression.
  \[
    B=\frac{\Delta^2}{2}(\mu_\alpha^2+\sigma_\alpha^2)\bigl(1-\mu_sE[1/s]\bigr)\leq0
    \ \text{unconditionally, yet the sign of }
    \left.\tfrac{\partial\operatorname{Var}(V(\theta))}{\partial\theta}\right|_{\theta=0}
    \ \text{is not.}
  \]

  \smallskip
  \textbf{2. Tested Definition 1's channel symmetry (\texttt{extensions/definition-asymmetry.md}).}
  Isolating the precision channel ($p$ only) gives a clean value gain
  $\theta/s$. Isolating the cost channel first looked \emph{impossible} under
  the paper's own normalization $c(0;\theta)=0$ --- I initially described this
  as failing ``eventually,'' which was imprecise: the contradiction is
  immediate, at every $e>0$. Relaxing that normalization (letting a fixed
  overhead shrink with $\theta$) produced an explicit, verified construction
  where the cost-only channel \emph{is} valid on the full domain.

  \begin{alertblock}{Takeaway}
    Both findings share a shape: a claim that looks unconditional on first
    read turns out to depend on an assumption the paper (or I) didn't
    surface --- condition (30) vs.\ the unconditional sign of $B$; the
    normalization $c(0;\theta)=0$ vs.\ Definition 1 itself.
  \end{alertblock}
\end{frame}

% 5 ------------------------------------------------------------------
\begin{frame}{Where I Did Not Believe the AI}
  \small
  \textbf{[TEMPLATE --- fill in after completing the hand derivation]}

  \begin{columns}[T,onlytextwidth]
    \column{0.72\textwidth}
      \textbf{What the model claimed:} \emph{[state the specific claim the
      LLM made that prompted the check --- quote or closely paraphrase it]}

      \smallskip
      \textbf{What I checked by hand:} \emph{[name the exact step: e.g.\ the
      Appendix A.1 ratio-monotonicity argument for Proposition 1, or a step
      from the extension's derivation]}

      \smallskip
      \textbf{My verdict:} \emph{[state one of: \textbf{correct} /
      \textbf{incorrect} / \textbf{unverifiable} --- and justify it in one or
      two sentences referencing what the hand derivation actually showed]}

    \column{0.24\textwidth}
      \centering
      \includegraphics[page=1,width=\linewidth,height=0.23\textheight,keepaspectratio]
        {hand/PLACEHOLDER.pdf}

      \scriptsize Handwritten derivation --- photo, crooked is fine
  \end{columns}

  \vspace{-0.3em}
  \begin{alertblock}{How to justify the verdict}
    \textbf{Correct:} the hand derivation reproduces the claim step-for-step.
    \textbf{Incorrect:} the hand derivation contradicts a specific step ---
    name it. \textbf{Unverifiable:} the claim depends on an assumption or
    computation that cannot be checked by hand in the time available --- say
    what would be needed.
  \end{alertblock}
\end{frame}

\end{document}
